% part: intuitionistic-logic
% chapter: propositions-as-types
% section: reduction

\documentclass[../../../include/open-logic-section]{subfiles}

\begin{document}

\olfileid{int}{pty}{red}

\olsection{راکمول}

د فطري استنتاج په !!{derivation}s کښې د
!!a{introduction} قاعدې پسې !!a{elimination} قاعده
داسې تاو راتاو دے چې «راکمېدل» ئې شوني دي۔ مثلاً
د $\Intro{\land}$ پسې $\Elim{\land}$ يو بل «لرې
کوي»، ځکه د $\Elim{\land}$ پايله تل د هغه او د
دواړو برخو يوه وي چې \Intro{\land} ئې جوړوي، او
د $\Intro{\land}$ دواړه مقدمې هماغه دواړه برخې
دي۔ نو که د يوې برخې !!a{derivation} پکار وي،
هماغه !!{derivation} نېغ کارولے شو، د او د
داخلولو او بيا ايستلو اړتيا نشته۔ داسې لرو:
\begin{gather*}
  \bottomAlignProof
  \AxiomC{}
  \RightLabel{$\delta_1$}
  \DeduceC{$!A_1$}
  \AxiomC{}
  \RightLabel{$\delta_2$}
  \DeduceC{$!A_2$}
  \RightLabel{$\Intro{\land}$}
  \BinaryInfC{$!A_1 \land !A_2$}
  \RightLabel{$\Elim{\land}_i$}
  \UnaryInfC{$!A_i$}
  \DisplayProof\bottomAlignProof
  \quad
  \redone
  \quad
  \AxiomC{}
  \RightLabel{$\delta_i$}
  \DeduceC{$!A_i$}
  \DisplayProof
\end{gather*}

د دې ساده کولو له مخې په اړوندو ثبوت ترمونو
کښې هم د راکمول ورته اړيکه اخيستلے شو۔ که $N_1$
د~$\delta_1$ ثبوت ترم وي او $N_2$ د~$\delta_2$
وي، نو ويلے شو چې د کيڼ ثبوت ثبوت ترم په يو ګام
کښې د ښي ثبوت ثبوت ترم ته راکمېږي:
\[
  \proj{i}{\pair{N_1}{N_2}} \redone N_i
\]
په ټايپ‌لرونکي لامبډا حساب کښې دا د ضرب ټايپ
\emph{د بېټا راکمول} قاعده ده۔
\textbf{د ژباړې د نسخې سمون (OLCMP-212):} د
جوړې ماکرو دوه آرګومېنټونه غواړي؛ د سرچينې دواړه
ترمونه په يو آرګومېنټ کښې وو۔ اوس په دوو جلا
آرګومېنټونو کښې دي، څو پروجکشن هماغه ټاکلے جز
راوباسي۔

په فطري استنتاج کښې د استنتاجونو داسې جوړه،
لکه د \Intro{\land} پسې \Elim{\land}، چې راکمېدل
پرې عملي کېږي، \emph{پرېکون} بلل کېږي۔ په
ټايپ‌لرونکي لامبډا حساب کښې د $\redone$ کيڼ ترم
\emph{رېډکس} او ښي ترم ئې \emph{راکم شوے ترم}
بلل کېږي۔ په بې‌ټايپه لامبډا حساب کښې يوازې
$(\lambd[x][N])Q$ رېډکس ګڼل کېږي؛ په ټايپ‌لرونکي
حساب کښې نحو داسې ترمونه هم رانغاړي چې
$\pair{N}{M}$، $\proj{i}{N}$، $\inj[!A]{i}{N}$،
$\dcase{N}{x_1}{M_1}{x_2}{M_2}$ او $\abort{!A}{N}$
پکې وي، او خپل اړوند رېډکسونه لري۔
\textbf{د ژباړې د نسخې سمون (OLCMP-213):} د يا
په جوړوونکي کښې بل ټايپ اختياري ټايپي نښه دے،
شاخص او ثبوت ترم ئې دوه اصلي آرګومېنټونه دي؛ د
سرچينې ناسم ماکرو ترتيب دلته د اصلي تعريف په
ترتيب برابر شو۔

په همدې ډول د يا دپاره راکمول لرو:
\begin{gather*}
  \bottomAlignProof
  \AxiomC{}
  \RightLabel{$\delta$}
  \DeduceC{$!A_i$}
  \RightLabel{$\Intro{\lor}$}
  \UnaryInfC{$!A_1 \lor !A_2$}
  \AxiomC{$\Discharge{!A_1}{x_1}$}
  \RightLabel{$\delta_1$}
  \DeduceC{$!C$}
  \AxiomC{$\Discharge{!A_2}{x_2}$}
  \RightLabel{$\delta_2$}
  \DeduceC{$!C$}
  \DischargeRule{$\Elim{\lor}$}{x_1,x_2}
  \TrinaryInfC{$!C$}
  \DisplayProof\bottomAlignProof
  \quad
  \redone
  \quad
  \AxiomC{}
  \RightLabel{$\delta$}
  \DeduceC{$!A_i$}
  \RightLabel{$\delta_i$}
  \DeduceC{$!C$}
  \DisplayProof
\end{gather*}
دې ته د ثبوت ترمونو لاندې راکمول برابر دي:
\begin{gather*}
  \dcase{\inj[!A_{3-i}]{i}{M}}{x_1}{N_1}{x_2}{N_2} \redone \Subst{N_i}{M}{x_i}
\end{gather*}
دلته $M$ د~$\delta$ ثبوت ترم او $N_i$ د~$\delta_i$
ثبوت ترم دے۔ دا د مجموعې ټايپونو د بېټا راکمول
قاعده ده۔ په عمومي نښه کښې $\Subst{M}{N_i}{x}$
دا معنا لري چې په $M$ کښې د متغير~$x$ ټول آزاد
موردونه په~$N_i$ بدل کړو۔
\textbf{د ژباړې د نسخې سمون (OLCMP-214):} د
ورګډ شوي ترم خپل ټايپ د مقدمې ټاکلې برخه ده؛ د
يا د جوړوونکي ټايپي نښه د بلې برخې ټايپ ښيي۔
نو په يوې/دوې څانګې کښې د بلې برخې ټايپ کارول
پکار دي؛ د سرچينې عين برخې نقل دلته برابر شو۔

د تابعې ټايپ راکمول د $\Intro\lif$ پسې د
$\Elim\lif$ د تاو راتاو لرې کولو سره برابر دي۔
\begin{gather*}
  \bottomAlignProof
  \AxiomC{$\Discharge{!A}{x}$}
  \RightLabel{$\delta$}
  \DeduceC{$!B$}
  \DischargeRule{$\Intro{\lif}$}{x}
  \UnaryInfC{$!A \lif !B$}
  \AxiomC{}
  \RightLabel{$\delta'$}
  \DeduceC{$!A$}
  \RightLabel{$\Elim{\lif}$}
  \BinaryInfC{$!B$}
  \DisplayProof\bottomAlignProof
  \quad
  \redone
  \quad
  \AxiomC{}
  \RightLabel{$\delta'$}
  \DeduceC{$!A$}
  \RightLabel{$\delta$}
  \DeduceC{$!B$}
  \DisplayProof
\end{gather*}
د ثبوت ترمونو دپاره دا هماغه عادي بېټا راکمول دي:
\begin{gather*}
  (\lambd[\typeof{x}{!A}][N]M)
  \redone \Subst{N}{M}{x}
\end{gather*}

په !!{proof}s کښې د راکمول د بدلونونو ترڅنګ
د ځاے بدلولو بدلونونه هم اخيستلے شو۔ دا پر هغو
(فرعي-)!!{proof}s عملي کېږي چې په \Elim{\lor}
پاې ته رسېږي او ورپسې د !!{elimination} بله
قاعده وي؛ مثلاً:
\begin{multline*}
  \bottomAlignProof
  \AxiomC{}
  \RightLabel{$\delta$}
  \DeduceC{$!A_i$}
  \RightLabel{$\Intro{\lor}$}
  \UnaryInfC{$!A_1 \lor !A_2$}
  \AxiomC{$\Discharge{!A_1}{x_1}$}
  \RightLabel{$\delta_1$}
  \DeduceC{$!C_1 \land !C_2$}
  \AxiomC{$\Discharge{!A_2}{x_2}$}
  \RightLabel{$\delta_2$}
  \DeduceC{$!C_1 \land !C_2$}
  \DischargeRule{$\Elim{\lor}$}{x_1,x_2}
  \TrinaryInfC{$!C_1 \land !C_2$}
  \RightLabel{\Elim{\land}[i]}
  \UnaryInfC{$!C_i$}
  \DisplayProof\bottomAlignProof
  \quad
  \redone
  \\
  \AxiomC{}
  \RightLabel{$\delta$}
  \DeduceC{$!A_i$}
  \RightLabel{$\Intro{\lor}$}
  \UnaryInfC{$!A_1 \lor !A_2$}
  \AxiomC{$\Discharge{!A_1}{x_1}$}
  \RightLabel{$\delta_1$}
  \DeduceC{$!C_1 \land !C_2$}
  \RightLabel{\Elim{\land}[i]}
  \UnaryInfC{$!C_i$}
  \AxiomC{$\Discharge{!A_2}{x_2}$}
  \RightLabel{$\delta_2$}
  \DeduceC{$!C_1 \land !C_2$}
  \RightLabel{\Elim{\land}[i]}
  \UnaryInfC{$!C_i$}
  \DischargeRule{$\Elim{\lor}$}{x_1,x_2}
  \TrinaryInfC{$!C_i$}
  \DisplayProof
\end{multline*}
که $M$ د هغه بشپړ يا-ثبوت ثبوت ترم وي چې
د~$\delta$ پسې د يا د داخلولو ګام هم پکې شامل
وي، او $N_1$ او $N_2$ په ترتيب د $\delta_1$ او
$\delta_2$ ثبوت ترمونه وي، نو د ثبوت ترمونو،
يعنې $\lambd$ ترمونو، اړوند راکمول قاعده دا ده:
\[
  \proj{i}{\dcase{M}{x_1}{N_1}{x_2}{N_2}} \redone 
\dcase{M}{x_1}{\proj{i}{N_1}}{x_2}{\proj{i}{N_2}}
\]
\textbf{د ژباړې د نسخې سمون (OLCMP-217):} په
رسم کښې د لومړني ثبوت پايله د يا يوه برخه ده،
خو د وروستي د يا د ايستلو ترم آرګومېنټ د يا بشپړ ثبوت
کوډ کوي۔ نو د هغه ترم ټاکنه د يا د داخلولو
ګام هم رانغاړي؛ اصلي رسمونه او د ځاے بدلولو
رسمي معادله عين دي۔

تاو راتاو يوازې د !!{proof}s له پاې څخه نۀ،
بلکې د !!{proof}s له منځه هم لرې کولے شو: په
داسې فرعي-!!{proof} د راکمول بدلون عملي کوو چې
په تاو راتاو پاې ته رسېږي۔ همدغه شان
$\beta$-راکمول د $\lambd$-ترمونو په فرعي ترمونو
هم عملي کېږي۔ که $N$ د $M$ فرعي ترم وي او
$N \redone N'$ وي، نو په~$M$ کښې فرعي ترم $N$
په $N'$ بدلولے شو، څو $M'$ ترلاسه کړو۔ په دې
حالت کښې $M \redone M'$ هم ليکو۔ که داسې لړۍ
$M = M_1 \redone M_2 \redone \dots \redone M_k = M'$
وي (د $k=1$، يعنې $M = M'$ حالت هم پکې شامل دے)،
نو $M \red M'$ ليکو۔
\textbf{د ژباړې د نسخې سمون (OLCMP-215):} د
سرچينې په لړۍ کښې منځنے ترم بې‌بدلونه يوګامي
ځان ته ورغلے و؛ دغه زيات نقل لرې شو۔ د صفرګامي
برابرۍ حالت د څوګامي اړيکې جلا شرط عين پاتې دے۔

هغه !!{proof} چې په هېڅ فرعي-!!{proof} کښې
ئې کوم بدلون نۀ شي عملي کېدے، \emph{عادي} بلل
کېږي۔ په همدې ډول هغه $\lambd$-ترم چې په هېڅ
فرعي ترم کښې ئې کوم بدلون نۀ شي عملي کېدے،
\emph{عادي} دے۔ عادي $\lambd$-ترم ته
\emph{عادي بڼه} هم وايي۔

\end{document}
